\section{Toposy elementarne i toposy Grothendiecka}\label{cat:sec:toposes}

Kategoria zbiorów ma nie tylko granice i obiekty wykładnicze, lecz także obiekt $\{0,1\}$, którego morfizmy klasyfikują podzbiory. W toposie zastąpimy te trzy własności warunkami uniwersalnymi. Punktem wyjścia są obiekty wykładnicze z rozdziału~\ref{sec:exponentials} oraz presnopy i snopy z rozdziału~\ref{cat:sec:grothendieck-topologies}.

\subsection{Klasyfikator podobiektów}

Niech $\mathcal E$ będzie lokalnie małą kategorią z granicami skończonymi. Podobiektem $X$ nazywamy klasę monomorfizmu $m:A\to X$ względem izomorfizmów nad $X$: morfizmy $m:A\to X$ oraz $n:B\to X$ reprezentują ten sam podobiekt, gdy istnieje izomorfizm $i:A\to B$ z $ni=m$. Piszemy $\operatorname{Sub}_{\mathcal E}(X)$ na zbiór tych klas, o ile jest zbiorem; w ogólnej lokalnie małej kategorii jest to a priori klasa.

\begin{definicja}[Klasyfikator podobiektów]\label{cat:def:subobject-classifier}
Klasyfikatorem podobiektów nazywamy obiekt $\Omega$ wraz z morfizmem $\mathrm{true}:1\to\Omega$ z obiektu końcowego, taki że dla każdego monomorfizmu $m:A\to X$ istnieje dokładnie jeden morfizm $\chi_m:X\to\Omega$, dla którego kwadrat
\begin{center}
\begin{tikzpicture}[>=stealth]
 \node (a) at (0,1.9) {$A$};
 \node (one) at (4.4,1.9) {$1$};
 \node (x) at (0,0) {$X$};
 \node (om) at (4.4,0) {$\Omega$};
 \draw[->] (a)--node[above] {$!_A$} (one);
 \draw[->] (a)--node[left] {$m$} (x);
 \draw[->] (one)--node[right] {$\mathrm{true}$} (om);
 \draw[->] (x)--node[below] {$\chi_m$} (om);
\end{tikzpicture}
\end{center}
jest kartezjański. Morfizm $\chi_m$ nazywamy morfizmem charakterystycznym podobiektu.
\end{definicja}

\begin{twierdzenie}\label{cat:thm:subobject-classifier-representation}
Klasyfikator podobiektów wyznacza naturalną bijekcję
\[
 \operatorname{Sub}_{\mathcal E}(X)\cong\Hom_{\mathcal E}(X,\Omega).
\]
W szczególności podobiekty każdego obiektu tworzą zbiór. Klasyfikator jest jednoznaczny z dokładnością do jedynego izomorfizmu zachowującego morfizm $\mathrm{true}$.
\end{twierdzenie}
\begin{proof}
Morfizm $\mathrm{true}$ jest monomorfizmem, gdyż dla dowolnego $T$ istnieje tylko jeden morfizm $T\to1$. Iloczyn włóknisty monomorfizmu jest monomorfizmem. Istotnie, dla monomorfizmu $n:B\to Y$, morfizmu $g:X\to Y$ i iloczynu włóknistego $P=X\times_Y B$ oznaczmy rzuty przez $p:P\to X$ i $q:P\to B$. Jeżeli $u,v:T\to P$ spełniają $pu=pv$, to $nqu=gpu=gpv=nqv$. Monomorficzność $n$ daje $qu=qv$, a jednoznaczność w definicji iloczynu włóknistego daje $u=v$. W szczególności projekcja $X\times_\Omega1\to X$ jest monomorfizmem.

Jeżeli dwa monomorfizmy są izomorficzne nad $X$, to kwadrat jednego jest kartezjański dokładnie wtedy, gdy jest nim kwadrat drugiego z tym samym dolnym morfizmem: izomorfizm przenosi w obie strony jedyne faktoryzacje. Stąd przyporządkowanie $[m]\mapsto\chi_m$ jest dobrze określone. W przeciwnym kierunku morfizm $f:X\to\Omega$ wyznacza klasę projekcji $X\times_\Omega1\to X$. Z definicji klasyfikatora morfizmem charakterystycznym tej projekcji jest $f$. Dla danego $m$ jego kwadrat jest kartezjański, więc $m$ jest izomorficzny nad $X$ z projekcją $X\times_\Omega1\to X$. Otrzymaliśmy wzajemnie odwrotne przyporządkowania. Ponieważ $\mathcal E$ jest lokalnie mała, $\Hom(X,\Omega)$ jest zbiorem; każdą klasę podobiektu można więc indeksować jednoznacznie elementem tego zbioru. W tym sensie kategoria jest dobrze uposażona w podobiekty (ang.\ \emph{well-powered}); nie twierdzimy, że wszystkie monomorfizmy do $X$ tworzą zbiór.

Dla $g:Y\to X$ przeciągamy $m$ do $Y$ przez iloczyn włóknisty. Morfizm charakterystyczny otrzymanego podobiektu to $\chi_m g$: istotnie, warunek faktoryzacji $h:T\to Y$ przez przeciągnięty podobiekt jest równoważny faktoryzacji $gh$ przez $m$, a ten jest równoważny równości $\chi_m gh=\mathrm{true}\,!_T$. Wszystkie faktoryzacje są jedyne, więc odpowiedni kwadrat jest kartezjański. Dowodzi to naturalności.

Dla dwóch klasyfikatorów $\Omega,\Omega'$ złożenie naturalnych bijekcji daje $\Hom(-,\Omega)\cong\Hom(-,\Omega')$. Z lematu Yonedy (twierdzenie~\ref{cat:thm:yoneda-contra}) transformacja ta ma postać $f\mapsto jf$, gdzie $j:\Omega\to\Omega'$ jest jej wartością na $\id_\Omega$. Transformacja odwrotna ma postać składania z pewnym $k':\Omega'\to\Omega$. Ich złożenia są identycznościami, więc podstawiając odpowiednie morfizmy identycznościowe, dostajemy $k'j=\id_\Omega$ i $jk'=\id_{\Omega'}$. Zatem $j$ jest izomorfizmem, a jego jednoznaczność wynika z wartości na $\id_\Omega$. W szczególności $j$ przenosi morfizm charakterystyczny $\id_1$ na morfizm charakterystyczny $\id_1$ w drugim klasyfikatorze, czyli $j\,\mathrm{true}=\mathrm{true}'$. Odwrotnie, jeśli izomorfizm $k$ spełnia $k\,\mathrm{true}=\mathrm{true}'$, to dla każdego $m$ kwadrat z dolnym morfizmem $k\chi_m$ jest kartezjański, ponieważ złożenie z izomorfizmem $k$ nie zmienia warunku równości w definicji iloczynu włóknistego. Stąd $k\chi_m=\chi'_m$ dla wszystkich $m$. Lemat Yonedy daje $k=j$.
\end{proof}

\subsection{Definicje i podstawowe przykłady}

\begin{definicja}[Topos elementarny]\label{cat:def:elementary-topos}
Toposem elementarnym nazywamy lokalnie małą kategorię $\mathcal E$, która ma wszystkie granice skończone, jest kartezjańsko domknięta (dla wszystkich $A,B$ istnieje obiekt wykładniczy $B^A$) oraz ma klasyfikator podobiektów.
\end{definicja}

\begin{definicja}[Topos Grothendiecka]\label{cat:def:grothendieck-topos}
Toposem Grothendiecka nazywamy kategorię równoważną kategorii $\operatorname{Sh}(\mathcal C,J)$ snopów zbiorów na pewnej małej kategorii $\mathcal C$ wyposażonej w topologię Grothendiecka $J$.
\end{definicja}

\begin{przyklad}[Kategoria zbiorów]\label{cat:ex:set-topos}
Kategoria $\Set$ jest toposem elementarnym. Ma granice skończone, a obiektem wykładniczym $B^A$ jest zbiór wszystkich funkcji $A\to B$ z ewaluacją. Klasyfikatorem jest $\Omega=\{0,1\}$ z $\mathrm{true}(*)=1$. Dla injekcji $m:A\to X$ definiujemy
\[
 \chi_m(x)=
 \begin{cases}1,&x\in m(A),\\0,&x\notin m(A).\end{cases}
\]
Jeżeli $h:T\to X$ ma $\chi_m h$ stale równą $1$, to dla każdego $t\in T$ istnieje jedyny $a_t\in A$ z $m(a_t)=h(t)$. Funkcja $t\mapsto a_t$ jest jedyną faktoryzacją $h$ przez $m$. Odwrotnie, dla każdego kartezjańskiego kwadratu tego typu, podstawiając $T=1$ i dowolny element $x\in X$, otrzymujemy $\chi_m(x)=1$ wtedy i tylko wtedy, gdy $x\in m(A)$. To dowodzi jednoznaczności. Argument obejmuje $A=\varnothing$ i $X=\varnothing$.

Kategoria $\Set$ jest również toposem Grothendiecka: na kategorii z jednym obiektem i tylko jego identycznością wybieramy topologię minimalną. Każdy presnop jest snopem, a wartość na jedynym obiekcie oraz składowa transformacji naturalnej dają izomorfizm kategorii snopów z $\Set$.
\end{przyklad}

\subsection{Topos presnopów}

\begin{lemat}\label{cat:lem:presheaf-monos-pointwise}
Morfizm $\alpha:F\to G$ presnopów na małej kategorii $\mathcal C$ jest monomorfizmem wtedy i tylko wtedy, gdy wszystkie funkcje $\alpha_U$ są injekcjami.
\end{lemat}
\begin{proof}
Jeżeli składowe są injekcjami, to dla $\alpha\beta=\alpha\gamma$ mamy $\alpha_U\beta_U=\alpha_U\gamma_U$, zatem $\beta_U=\gamma_U$ dla wszystkich $U$. Odwrotnie, jeśli $\alpha$ jest monomorfizmem i $s,t\in F(U)$ mają ten sam obraz, lemat Yonedy odpowiada im transformacjami $h_U\to F$. Ich złożenia z $\alpha$ są równe, gdyż odpowiadają tej samej sekcji $\alpha_U(s)=\alpha_U(t)$ w $G(U)$. Monomorficzność daje równość transformacji, a bijekcja Yonedy daje $s=t$.
\end{proof}

Wobec lematu każdy podobiekt presnopu można reprezentować podfunktorem: dla $\alpha:F\to G$ bierzemy $A(U)=\alpha_U(F(U))$. Naturalność $\alpha$ gwarantuje, że ograniczenia $G(f)$ przeprowadzają $A(U)$ do $A(V)$. Funkcje $\bar\alpha_U:F(U)\to A(U)$ otrzymane przez ograniczenie przeciwdziedziny $\alpha_U$ są bijekcjami. Dla $f:V\to U$ naturalność daje $A(f)\bar\alpha_U=\bar\alpha_V F(f)$, a więc $\bar\alpha_V^{-1}A(f)=F(f)\bar\alpha_U^{-1}$. Odwrotności tworzą transformację naturalną. Wobec tego $\alpha$ i włączenie $A\hookrightarrow G$ reprezentują ten sam podobiekt. Jeżeli dwa podfunktory $A,B\subseteq G$ reprezentują ten sam podobiekt, to izomorfizm nad $G$ identyfikuje każdy element $A(U)$ z tym samym elementem w $B(U)$ i odwrotnie. Stąd $A(U)=B(U)$ dla każdego $U$, a ograniczenia są odziedziczone z $G$, więc $A=B$.

\begin{twierdzenie}[Klasyfikator w kategorii presnopów]\label{cat:thm:presheaf-classifier}
W $\widehat{\mathcal C}$ klasyfikatorem podobiektów jest presnop
\[
 \Omega(U)=\{R:R\text{ jest sitem na }U\},\qquad
 \Omega(f)(R)=f^*R.
\]
Morfizm $\mathrm{true}:1\to\Omega$ ma jako wartość na $U$ sito maksymalne $h_U$.
\end{twierdzenie}
\begin{proof}
Zbiór sit istnieje z małości $\mathcal C$. Równości $(fg)^*R=g^*f^*R$ i $\id_U^*R=R$ ze stwierdzenia~\ref{cat:prop:sieve-pullback} dowodzą, że $\Omega$ jest presnopem. Przeciągnięcie sita maksymalnego jest maksymalne, więc $\mathrm{true}$ jest naturalne.

Dla podfunktora $A\subseteq F$ i $s\in F(U)$ kładziemy
\[
 \chi_{A,U}(s)=\{f:V\to U:s|_f\in A(V)\}.
\]
Jeżeli $f$ należy do tego zbioru i $g:W\to V$, to $s|_{fg}=(s|_f)|_g\in A(W)$, ponieważ $A$ jest podfunktorem. Otrzymujemy sito. Dla $a:V\to U$ i $g:W\to V$ zachodzi
\[
 g\in\chi_{A,V}(s|_a)
 \iff s|_{ag}\in A(W)
 \iff g\in a^*\chi_{A,U}(s),
\]
więc $\chi_A$ jest naturalne. Ponadto $\chi_{A,U}(s)=h_U$ dokładnie wtedy, gdy $\id_U$ należy do tego sita, czyli $s\in A(U)$. W każdej składowej przeciwobrazem jednoelementowego zbioru $\{h_U\}$ przez $\chi_{A,U}$ jest dokładnie $A(U)$. Kwadrat składowych jest zatem iloczynem włóknistym w $\Set$: funkcja do $F(U)$ faktoryzuje się jednoznacznie przez $A(U)$ dokładnie wtedy, gdy jej obraz jest zawarty w $A(U)$. Z punktowego obliczania granic presnopów (stwierdzenie~\ref{cat:prop:presheaf-representables}) wynika, że kwadrat klasyfikujący $A$ jest kartezjański.

Niech $\psi:F\to\Omega$ będzie innym morfizmem, którego kwadrat jest kartezjański. Punktowo mamy $\psi_V(t)=h_V$ dokładnie wtedy, gdy $t\in A(V)$. Dla $f:V\to U$ naturalność oraz stwierdzenie~\ref{cat:prop:sieve-pullback} dają
\[
 f\in\psi_U(s)
 \iff f^*\psi_U(s)=h_V
 \iff \psi_V(s|_f)=h_V
 \iff s|_f\in A(V).
\]
Zatem $\psi=\chi_A$. Każdy monomorfizm reprezentuje podfunktor na mocy poprzedniego lematu, co kończy dowód.
\end{proof}

\begin{twierdzenie}[Obiekty wykładnicze presnopów]\label{cat:thm:presheaf-exponential}
Dla presnopów $F,G$ istnieje obiekt wykładniczy $G^F$, dany przez
\[
 (G^F)(U)=\Nat(h_U\times F,G).
\]
Dla $a:V\to U$ ograniczenie $\theta\in (G^F)(U)$ jest równe $\theta\circ(h_a\times\id_F)$, gdzie $h_a:h_V\to h_U$ jest obrazem $a$ przez zanurzenie Yonedy.
\end{twierdzenie}
\begin{proof}
Małość $\mathcal C$ zapewnia, że $\Nat(h_U\times F,G)$ jest zbiorem. Złożenie transformacji naturalnych jest naturalne, więc podana formuła określa funkcję ograniczenia. Dla $a:V\to U$ i $b:W\to V$ mamy
\[
 (\theta|_a)|_b=\theta(h_a\times\id_F)(h_b\times\id_F)
 =\theta(h_{ab}\times\id_F)=\theta|_{ab},
\]
gdyż $h_{ab}=h_ah_b$ i iloczyn morfizmów zachowuje składanie. Z $h_{\id_U}=\id_{h_U}$ wynika $\theta|_{\id_U}=\theta$. Otrzymujemy presnop. Ewaluację definiujemy punktowo przez
\[
 \ev_U(\theta,s)=\theta_U(\id_U,s).
\]
Dla $a:V\to U$ naturalność $\theta$ daje
\[
 \ev_U(\theta,s)|_a=\theta_V(a,s|_a)
 =\ev_V(\theta|_a,s|_a),
\]
więc $\ev:G^F\times F\to G$ jest transformacją naturalną.

Dla dowolnego presnopu $H$ i morfizmu $\beta:H\times F\to G$ konstruujemy $\widetilde\beta:H\to G^F$. Dla $t\in H(U)$ i $f:V\to U$, $s\in F(V)$ definiujemy
\[
 (\widetilde\beta_U(t))_V(f,s)=\beta_V(t|_f,s).
\]
Dla $g:W\to V$ ograniczenie tej wartości jest, z naturalności $\beta$, równe $\beta_W(t|_{fg},s|_g)$, co dowodzi naturalności w $V$. Dla $a:V\to U$ obie strony równości $\widetilde\beta_V(t|_a)=\widetilde\beta_U(t)|_a$ przyjmują na $(g:W\to V,s)$ wartość $\beta_W(t|_{ag},s)$. Zatem $\widetilde\beta$ jest morfizmem presnopów. Wprost z definicji
\[
 \ev_U(\widetilde\beta_U(t),s)=\beta_U(t,s).
\]

Jeśli $\gamma:H\to G^F$ spełnia $\ev(\gamma\times\id_F)=\beta$, to dla $f:V\to U$ i $s\in F(V)$ naturalność $\gamma$ daje
\[
 (\gamma_U(t))_V(f,s)
 =(\gamma_U(t)|_f)_V(\id_V,s)
 =\ev_V(\gamma_V(t|_f),s)
 =\beta_V(t|_f,s).
\]
Stąd $\gamma=\widetilde\beta$, ponieważ składowe transformacji są funkcjami i zgadzają się na każdym argumencie. Bijekcja pomiędzy morfizmami $H\times F\to G$ i $H\to G^F$ jest naturalna w $H$: dla $\delta:K\to H$ oba morfizmy $\widetilde{\beta(\delta\times\id_F)}$ i $\widetilde\beta\,\delta$ po złożeniu z ewaluacją dają $\beta(\delta\times\id_F)$, więc są równe z właśnie udowodnionej jednoznaczności. Otrzymaliśmy wymaganą własność uniwersalną. W przypadku pustych wartości presnopów wzory nadal definiują funkcje, a warunki odnoszące się do nieistniejących elementów są puste.
\end{proof}

\begin{wniosek}\label{cat:cor:presheaf-topos}
Dla każdej małej kategorii $\mathcal C$ kategoria $\widehat{\mathcal C}$ jest toposem elementarnym oraz toposem Grothendiecka.
\end{wniosek}
\begin{proof}
Lokalna małość i istnienie granic skończonych wynikają z definicji presnopów i stwierdzenia~\ref{cat:prop:presheaf-representables}. Twierdzenia~\ref{cat:thm:presheaf-classifier} i~\ref{cat:thm:presheaf-exponential} dają pozostałe aksjomaty toposu elementarnego. Przy topologii minimalnej wszystkie presnopy są snopami (rozdział~\ref{cat:sec:grothendieck-topologies}), więc $\widehat{\mathcal C}$ jest również toposem Grothendiecka.
\end{proof}

\begin{przyklad}[Wartości prawdy w kategorii presnopów]\label{cat:ex:presheaf-truth-values}
Niech $\mathcal C$ ma dwa obiekty $V,U$ i jedyny nieidentycznościowy morfizm $f:V\to U$. Na $V$ są dwa sita, $\varnothing$ i $h_V$, zaś na $U$ trzy: $\varnothing,\{f\},h_U$. Ograniczenie $\Omega(f)$ wysyła sito puste na puste, a pozostałe dwa na $h_V$. W szczególności istnieje lokalna wartość prawdy na $U$, która nie jest maksymalna, lecz staje się maksymalna po ograniczeniu do $V$.

Dla presnopu $F=1$ wybierzmy podfunktor $A$ z $A(U)=\varnothing$ i $A(V)=\{*\}$. Ograniczenie z pustego zbioru do zbioru jednoelementowego istnieje, więc jest to rzeczywiście podfunktor. Jego morfizm charakterystyczny przyjmuje na $U$ wartość $\{f\}$, ponieważ jedyna sekcja $F(U)$ nie należy do $A(U)$, ale jej ograniczenie należy do $A(V)$. Na $V$ wartość jest maksymalna. Tak otrzymujemy konkretny przykład wartości prawdy różnej od klasycznych dwóch wartości w $\Set$.
\end{przyklad}


\subsection{Sita domknięte względem topologii Grothendiecka}

Ustalmy małe miejsce $(\mathcal C,J)$. Nie zakładamy subkanoniczności: presnopy $h_U$ nie muszą być snopami.

\begin{definicja}[Domknięcie sita]\label{cat:def:j-closed-sieve}
Dla sita $R$ na $U$ kładziemy
\[
 c_J(R)=\{f:V\to U:f^*R\in J(V)\}.
\]
Sito $R$ nazywamy $J$-domkniętym, jeżeli $c_J(R)=R$.
\end{definicja}

\begin{lemat}\label{cat:lem:j-sieve-closure}
Operacja $c_J$ przyjmuje wartości w sitach, jest ekstensywna, monotoniczna i idempotentna, zachowuje skończone przecięcia oraz spełnia $a^*c_J(R)=c_J(a^*R)$.
\end{lemat}
\begin{proof}
Dla $f\in c_J(R)$ i $g:W\to V=\operatorname{dom}f$ sito $(fg)^*R=g^*f^*R$ pokrywa $W$ przez stabilność. Zatem $c_J(R)$ jest sitem. Jeżeli $f\in R$, to $f^*R=h_V$ pokrywa $V$, więc $f\in c_J(R)$. Jeżeli $R\subseteq S$, to $f^*R\subseteq f^*S$; z lematu~\ref{cat:lem:covering-refinement} wynika monotoniczność $c_J$.

Dla $a:V\to U$ i $g:W\to V$ mamy
\[
 g\in a^*c_J(R)\iff (ag)^*R\in J(W)
 \iff g^*(a^*R)\in J(W)\iff g\in c_J(a^*R).
\]
Jeżeli $f:V\to U$ należy do $c_J(c_J(R))$, to $f^*c_J(R)=c_J(f^*R)$ pokrywa $V$. Dla każdego $g:W\to V$ w tym sicie $g^*f^*R$ pokrywa $W$. Aksjomat lokalny daje $f^*R\in J(V)$, czyli $f\in c_J(R)$. Inkluzja odwrotna wynika z ekstensywności, co dowodzi idempotentności.

Monotoniczność daje $c_J(R\cap S)\subseteq c_J(R)\cap c_J(S)$. Jeżeli $f$ należy do prawej strony, to $f^*R,f^*S$ są pokrywające, a z lematu~\ref{cat:lem:covering-refinement} także $f^*(R\cap S)=f^*R\cap f^*S$ pokrywa. Stąd równość. Sito maksymalne pozostaje maksymalne po domknięciu, więc zachodzi również przypadek pustego przecięcia.
\end{proof}

\subsection{Klasyfikator podobiektów snopów}

\begin{twierdzenie}\label{cat:thm:closed-sieves-sheaf}
Presnop
\[
 \Omega_J(U)=\{R:R\text{ jest }J\text{-domkniętym sitem na }U\},
 \qquad \Omega_J(f)(R)=f^*R
\]
jest snopem.
\end{twierdzenie}
\begin{proof}
Lemat~\ref{cat:lem:j-sieve-closure} zapewnia, że ograniczenia przyjmują wartości w sitach domkniętych, a stwierdzenie~\ref{cat:prop:sieve-pullback} daje funktorialność. Niech $P\in J(U)$ i $(R_f)_{f\in P}$ będzie zgodną rodziną, czyli $R_{fg}=g^*R_f$.

Definiujemy $S$ na $U$: morfizm $h:W\to U$ należy do $S$ wtedy i tylko wtedy, gdy dla każdego $g:Z\to W$ takiego, że $hg\in P$, mamy $\id_Z\in R_{hg}$. Jeżeli $h\in S$ i $k:T\to W$, to warunek dla $hk$ i dowolnego $g:Z\to T$ z $hkg\in P$ wynika z warunku dla $h$ i $kg$. Zatem $S$ jest sitem.

Niech $h^*S$ pokrywa $W$ i niech $g:Z\to W$ spełnia $hg\in P$. Dla $k:T\to Z$ w pokrywającym sicie $g^*h^*S$ mamy $hgk\in S\cap P$, więc z definicji $S$ (dla identyczności $\id_T$) wynika $\id_T\in R_{hgk}=k^*R_{hg}$. Wobec tego $k\in R_{hg}$. Sito $R_{hg}$ zawiera pokrywające $g^*h^*S$, zatem pokrywa $Z$. Jego domkniętość daje $\id_Z\in R_{hg}$. Ponieważ $g$ było dowolne, $h\in S$. Otrzymaliśmy $c_J(S)\subseteq S$, a inkluzja $S\subseteq c_J(S)$ wynika z ekstensywności udowodnionej w lemacie~\ref{cat:lem:j-sieve-closure}. Stąd $c_J(S)=S$.

Dla $f:V\to U$ w $P$ oraz $g:W\to V$ warunek $g\in f^*S$ oznacza, że $\id_Z\in R_{fgk}=(gk)^*R_f$ dla każdego $k:Z\to W$. Jest to równoważne $g\in R_f$: w jedną stronę korzystamy z domknięcia sita względem składania z prawej, w drugą podstawiamy $k=\id_W$. Otrzymujemy $f^*S=R_f$.

Jeżeli $S'$ jest innym domkniętym sitem o tych samych ograniczeniach, to dla $h\in S'$ i $hg\in P$ mamy $hg\in S'$, czyli $\id_{\operatorname{dom}g}\in (hg)^*S'=R_{hg}$. Stąd $S'\subseteq S$. Dla $h\in S$ definicja daje $h^*P\subseteq h^*S'$. Sito $h^*P$ pokrywa, więc $h^*S'$ pokrywa, a domkniętość $S'$ daje $h\in S'$. Otrzymujemy $S=S'$. W przypadku pustego sita pokrywającego $P$ warunek definiujący $S$ jest pusty i $S=h_U$; powyższy argument nadal zapewnia jednoznaczność.
\end{proof}

\begin{lemat}\label{cat:lem:sheaf-monos-pointwise}
Monomorfizmy snopów są dokładnie morfizmami punktowo injektywnymi.
\end{lemat}
\begin{proof}
Punktowa injektywność daje monomorficzność w presnopach na mocy lematu~\ref{cat:lem:presheaf-monos-pointwise}, a więc także w snopach. Odwrotnie, sprzężenie snopifikacji i lemat Yonedy dają naturalne bijekcje
\[
 \Hom_{\operatorname{Sh}(\mathcal C,J)}(a h_U,F)
 \cong\Nat(h_U,F)\cong F(U).
\]
Jeżeli monomorfizm $\alpha:F\to G$ spełnia $\alpha_U(s)=\alpha_U(t)$, to odpowiadające sekcjom morfizmy $a h_U\to F$ mają równe złożenia z $\alpha$ dzięki naturalności bijekcji. Monomorficzność daje równość tych morfizmów, a więc $s=t$. Nie potrzeba zakładać, że $h_U$ jest snopem.
\end{proof}

\begin{twierdzenie}[Klasyfikator w kategorii snopów]\label{cat:thm:sheaf-classifier}
Klasyfikatorem podobiektów snopów jest $\Omega_J$ wraz z morfizmem $\mathrm{true}:1\to\Omega_J$ wybierającym sito maksymalne na każdym obiekcie.
\end{twierdzenie}
\begin{proof}
Sita maksymalne są domknięte i ich przeciągnięcia są maksymalne, więc $\mathrm{true}$ jest naturalne. Obiekt końcowy jest snopem z twierdzenia~\ref{cat:thm:sheaf-limits}.

Na mocy poprzedniego lematu każdy monomorfizm snopów reprezentujemy jako włączenie podfunktora $A\subseteq F$ izomorficznego z jego dziedziną, a zatem będącego snopem. Z twierdzenia~\ref{cat:thm:presheaf-classifier} morfizm charakterystyczny presnopów przypisuje $s\in F(U)$ sito
\[
 R_s=\{f:V\to U:s|_f\in A(V)\}.
\]
Jeżeli $h:W\to U$ spełnia $h^*R_s\in J(W)$, to sekcje $s|_{hg}$ dla $g:Z\to W$ w $h^*R_s$ należą do $A(Z)$ i tworzą zgodną rodzinę w $A$, ponieważ są ograniczeniami jednej sekcji w $F$, a $A$ jest podfunktorem. Sklejenie w snopie $A$ daje $t\in A(W)$. W $F(W)$ sekcje $t$ i $s|_h$ mają równe ograniczenia na pokrywającym sicie $h^*R_s$, więc z rozdzieloności $F$ są równe. Stąd $h\in R_s$, a z ekstensywności $c_J(R_s)=R_s$.

Morfizm charakterystyczny presnopów przyjmuje zatem wartości w $\Omega_J$ i określa morfizm snopów. Jego kwadrat klasyfikujący jest punktowo kartezjański, gdyż $R_s=h_U$ dokładnie wtedy, gdy $s\in A(U)$. Granice snopów obliczamy punktowo (twierdzenie~\ref{cat:thm:sheaf-limits}), więc kwadrat jest kartezjański w snopach. Jeżeli inny morfizm do $\Omega_J$ klasyfikuje $A$, to po złożeniu z włączeniem $\Omega_J\to\Omega$ klasyfikuje $A$ w presnopach: punktowo przeciwobraz sita maksymalnego pozostaje ten sam. Jednoznaczność klasyfikacji presnopów i injektywność włączenia dają równość obu morfizmów.
\end{proof}

\subsection{Obiekty wykładnicze snopów}

\begin{twierdzenie}\label{cat:thm:sheaf-exponential}
Jeżeli $F,G$ są snopami, to ich obiekt wykładniczy $G^F$ w presnopach jest snopem i jest obiektem wykładniczym w kategorii snopów. Ponadto naturalnie
\[
 (G^F)(U)\cong
 \Hom_{\operatorname{Sh}(\mathcal C,J)}(a h_U\times F,G).
\]
\end{twierdzenie}
\begin{proof}
Dla $R\in J(U)$ włączenie $i:R\to h_U$ staje się izomorfizmem po snopifikacji. Istotnie, dla $t:V\to U$ w $h_U(V)$ sito $t^*R$ pokrywa $V$, a dla każdego $g:W\to V$ w tym sicie morfizm $tg$ należy do $R(W)$ i jest przeciwobrazem ograniczenia $t|_g$. Włączenie $i$ jest punktowo injektywne, więc dla równych obrazów sekcje są już równe, a jako sito pokrywające można wybrać maksymalne. Spełnione są oba warunki twierdzenia~\ref{cat:thm:local-sheafification-isomorphism}, zatem $ai$ jest izomorfizmem.

Dla dowolnego presnopu $P$ mamy naturalne bijekcje
\[
 \begin{aligned}
 \Nat(P,G^F)&\cong\Nat(P\times F,G)\\
 &\cong\Hom_{\operatorname{Sh}(\mathcal C,J)}(a(P\times F),G)\\
 &\cong\Hom_{\operatorname{Sh}(\mathcal C,J)}(aP\times F,G).
 \end{aligned}
\]
Pierwsza to własność wykładnicza presnopów, druga to sprzężenie snopifikacji (ponieważ $G$ jest snopem), trzecia wynika z zachowywania iloczynów przez $a$ (twierdzenie~\ref{cat:thm:sheafification-left-exact}) i izomorfizmu $aF\cong F$ (wniosek~\ref{cat:cor:sheafification-reflective}). Przez naturalność w $P$ ograniczenie $\Nat(h_U,G^F)\to\Nat(R,G^F)$ odpowiada składaniu z $ai\times\id_F$. Jest to izomorfizm o odwrotności $(ai)^{-1}\times\id_F$, więc ograniczenie jest bijekcją. Warunek snopu w języku Yonedy (definicja~\ref{cat:def:site-sheaf}) daje tezę.

Dla snopu $H$ iloczyn $H\times F$ jest snopem, ponieważ granice są punktowe. Pełność włączenia snopów i wykładnicza bijekcja presnopów dają zatem wykładniczą bijekcję w kategorii snopów. Ostatni wzór wynika z podstawienia $P=h_U$ w ciągu powyższych bijekcji i lematu Yonedy.
\end{proof}

\begin{twierdzenie}\label{cat:thm:grothendieck-is-elementary}
Każdy topos Grothendiecka jest toposem elementarnym.
\end{twierdzenie}
\begin{proof}
Dla małego miejsca kategoria snopów jest lokalnie mała jako pełna podkategoria lokalnie małej kategorii presnopów. Ma granice skończone z twierdzenia~\ref{cat:thm:sheaf-limits}, klasyfikator podobiektów z twierdzenia~\ref{cat:thm:sheaf-classifier} oraz obiekty wykładnicze z twierdzenia~\ref{cat:thm:sheaf-exponential}. Spełnia więc definicję~\ref{cat:def:elementary-topos}. Równoważność kategorii przenosi te własności: bijekcje zbiorów morfizmów oraz istotna surjektywność przenoszą warunki uniwersalne granic i obiektów wykładniczych; pełność i wierność wraz z istotną surjektywnością zachowują i odzwierciedlają monomorfizmy. Własność uniwersalna klasyfikatora, wyrażona przez monomorfizmy i kwadraty kartezjańskie, również się przenosi. Lokalna małość przenosi się przez bijekcje zbiorów morfizmów.
\end{proof}

\begin{przyklad}[Klasyfikator na przestrzeni topologicznej]\label{cat:ex:topological-truth-values}
Dla miejsca $\mathcal O(X)$ z topologią pokryć otwartych mamy
\[
 \Omega_J(U)\cong\{W\in\mathcal O(X):W\subseteq U\}.
\]
Dla sita $R$ na $U$ niech $W=\bigcup_{V\in R}V$. Przeciągnięcie sita do otwartego $T\subseteq U$ składa się z otwartych podzbiorów $T$ należących do $R$. Ich suma jest równa $T\cap W$: zawiera każdy $T\cap V$ dla $V\in R$, a każdy składnik sumy jest zawarty w $T\cap W$. Przeciągnięcie pokrywa $T$ dokładnie wtedy, gdy $T\subseteq W$. Zatem $c_J(R)$ to sito wszystkich otwartych podzbiorów $W$. Sito tej postaci jest domknięte, gdyż suma jego elementów jest równa $W$ (sam $W$ do niego należy). Dwa takie sita są równe tylko wtedy, gdy ich maksymalne elementy są równe. Stąd bijekcja. Ograniczenie do $T$ odpowiada przecięciu $W\cap T$, a $\mathrm{true}$ odpowiada całemu $U$. Dla pustego $U$ otrzymujemy jedną wartość, zgodnie z aksjomatem snopu.
\end{przyklad}


\subsection{Obiekty potęgowe i logika podobiektów}

\begin{twierdzenie}[Obiekt potęgowy]\label{cat:thm:power-object}
W toposie elementarnym dla każdego obiektu $A$ istnieje obiekt $\mathcal P(A)=\Omega^A$ reprezentujący funktor $T\mapsto\operatorname{Sub}_{\mathcal E}(T\times A)$. Istnieje uniwersalny podobiekt $M_A\hookrightarrow\mathcal P(A)\times A$, nazywany relacją przynależności.
\end{twierdzenie}
\begin{proof}
Klasyfikator i własność wykładnicza dają naturalne bijekcje
\[
 \operatorname{Sub}_{\mathcal E}(T\times A)
 \cong\Hom(T\times A,\Omega)
 \cong\Hom(T,\Omega^A).
\]
Niech $M_A$ będzie podobiektem klasyfikowanym przez ewaluację $\ev:\Omega^A\times A\to\Omega$. Dla podobiektu $R\hookrightarrow T\times A$ jego morfizm charakterystyczny $\chi_R$ ma jedyny odpowiednik $r:T\to\Omega^A$ spełniający $\chi_R=\ev(r\times\id_A)$. Z naturalności klasyfikacji (twierdzenie~\ref{cat:thm:subobject-classifier-representation}) wynika, że przeciągnięcie $M_A$ przez $r\times\id_A$ jest klasyfikowane przez $\chi_R$, a zatem jest równe $R$ jako podobiekt. Jednoznaczność $r$ wynika z bijekcji wykładniczej.
\end{proof}

Porządkujemy podobiekty obiektu $X$ przez faktoryzację: dla $m:A\to X$ i $n:B\to X$ przyjmujemy $[m]\leq[n]$, gdy istnieje $u:A\to B$ z $nu=m$. Jeżeli zastępujemy $m,n$ przez $m':A'\to X,n':B'\to X$ z izomorfizmami $i:A\to A'$ i $j:B\to B'$ spełniającymi $m'i=m$, $n'j=n$, to $n'(jui^{-1})=m'$. Odwrotną implikację otrzymujemy za pomocą $i^{-1},j^{-1}$. Porządek nie zależy zatem od reprezentantów. Morfizm $u$ jest jedyny z monomorficzności $n$; jeśli $us=ut$, to $ms=mt$, więc $s=t$. Zwrotność wynika z faktoryzacji przez identyczność, a przechodniość ze składania faktoryzacji. Jeżeli ponadto $mv=n$ dla $v:B\to A$, to $nuv=mv=n$ i $mvu=nu=m$, skąd $uv=\id_B$ i $vu=\id_A$. Otrzymujemy antysymetrię na klasach podobiektów.

\begin{definicja}[Algebra Heytinga]\label{cat:def:heyting-algebra}
Algebrą Heytinga nazywamy ograniczoną kratę $(L,\wedge,\vee,0,1)$ z operacją implikacji $\Rightarrow$ spełniającą
\[
 c\leq(a\Rightarrow b)\quad\Longleftrightarrow\quad c\wedge a\leq b.
\]
Negację definiujemy jako $\neg a=(a\Rightarrow0)$.
\end{definicja}

Implikacja jest jednoznacznie określona przez warunek uniwersalny: jeżeli dwa elementy spełniają go dla wszystkich $c$, to każdy z nich jest nie większy od drugiego. Krata taka jest rozdzielna, ponieważ dla ustalonego $a$ operacja $c\mapsto c\wedge a$ jest lewym sprzężeniem operacji $b\mapsto a\Rightarrow b$ w kategoriach uporządkowanych. Bezpośrednio, dla dowolnego $d$,
\[
 (x\vee y)\wedge a\leq d
 \iff x\vee y\leq(a\Rightarrow d)
 \iff (x\wedge a)\vee(y\wedge a)\leq d.
\]
Podstawiając jako $d$ kolejno oba porównywane wyrażenia, dostajemy rozdzielność.

\begin{twierdzenie}[Algebra Heytinga podobiektów w toposie elementarnym]\label{cat:thm:elementary-heyting}
Dla każdego obiektu $X$ toposu elementarnego $\operatorname{Sub}_{\mathcal E}(X)$ jest algebrą Heytinga. Przeciąganie podobiektów przez dowolny morfizm zachowuje $0,1,\wedge,\vee,\Rightarrow$.
\end{twierdzenie}
\begin{proof}
Elementem największym jest $X\hookrightarrow X$, a przecięcie dwóch podobiektów jest ich iloczynem włóknistym nad $X$: faktoryzacja przez niego jest równoważna faktoryzacji przez oba podobiekty. Przeciągnięcie identyczności obiektu $X$ przez $g:Y\to X$ daje identyczność $Y$. Dla podobiektów $A,B$ obiektu $X$ i morfizmu $h:T\to Y$ faktoryzacja $h$ przez $g^*(A\wedge B)$ jest równoważna faktoryzacji $gh$ przez oba podobiekty $A,B$. To z kolei jest równoważne faktoryzacji $h$ przez oba podobiekty $g^*A,g^*B$, czyli przez ich przecięcie. Wszystkie faktoryzacje są jedyne z monomorficzności włączeń. Podstawiając za $h$ włączenia obu porównywanych podobiektów $Y$, dostajemy obie nierówności, a zatem $g^*(A\wedge B)=g^*A\wedge g^*B$ w porządku podobiektów.

Skonstruujemy najpierw implikację, nie zakładając jeszcze istnienia elementu najmniejszego ani sum. Niech $\mu:\Omega\times\Omega\to\Omega$ klasyfikuje iloczyn monomorfizmu $\mathrm{true}:1\to\Omega$ z nim samym (identyfikujemy $1\times1$ z $1$). Dla $a,b:T\to\Omega$ morfizm $\mu(a,b)$ klasyfikuje przecięcie odpowiadających im podobiektów: przeciągnięcie iloczynu $\mathrm{true}\times\mathrm{true}$ przez $(a,b)$ daje ich iloczyn włóknisty nad $T$.

Niech $i:L\hookrightarrow\Omega\times\Omega$ będzie equalizatorem $\mu$ i pierwszego rzutu. Dla pary $(a,b):T\to\Omega\times\Omega$ faktoryzacja przez $L$ oznacza równość $\mu(a,b)=a$, czyli inkluzję podobiektu klasyfikowanego przez $a$ w podobiekcie klasyfikowanym przez $b$. Niech $\lambda:\Omega\times\Omega\to\Omega$ klasyfikuje $i$.

Dla podobiektów $A,B\hookrightarrow X$ o morfizmach charakterystycznych $a,b$ definiujemy $D$ jako podobiekt klasyfikowany przez $\lambda(a,b)$. Dla $C\hookrightarrow X$ o włączeniu $m_C$ zachodzi
\[
 \begin{aligned}
 C\leq D
 &\iff (am_C,bm_C)\text{ faktoryzuje się przez }L\\
 &\iff m_C^*A\leq m_C^*B\\
 &\iff C\wedge A\leq B.
 \end{aligned}
\]
Pierwszy warunek wynika z kwadratu klasyfikującego $L$. Ostatnia równoważność wymaga utożsamienia podobiektów $C$ z podobiektami $X$ zawartymi w $C$. Składanie monomorfizmu $n:E\to C$ z $m_C$ daje monomorfizm $m_Cn$, a każdy podobiekt zawarty w $C$ ma taki reprezentant dzięki definicji porządku. Dla $n,n'$ warunek $m_Cn'u=m_Cn$ jest równoważny $n'u=n$, z monomorficzności $m_C$. Składanie zachowuje więc i odzwierciedla porządek. Przeciągnięcia $m_C^*A,m_C^*B$ odpowiadają odpowiednio $C\wedge A,C\wedge B$, ponieważ ich złożone włączenia są projekcjami iloczynów włóknistych nad $X$ do $X$. Warunek $C\wedge A\leq C\wedge B$ jest równoważny $C\wedge A\leq B$, gdyż $C\wedge A\leq C$. Tak skonstruowane $D$ jest implikacją $A\Rightarrow B$. Dla $g:Y\to X$ naturalność klasyfikacji daje morfizm charakterystyczny $g^*D$ równy $\lambda(a,b)g=\lambda(ag,bg)$, co dowodzi stabilności implikacji przy przeciąganiu.

Do konstrukcji elementu najmniejszego i sumy użyjemy kwantyfikacji po obiekcie $\Omega$. Dla dowolnego obiektu $K$ i podobiektu $R\hookrightarrow X\times K$ o morfizmie charakterystycznym $r$ niech $\forall_K R$ będzie equalizatorem dwóch morfizmów $X\to\Omega^K$: transpozycji wykładniczej $r$ oraz morfizmu stałego o wartości odpowiadającej $\mathrm{true}\,!_K:K\to\Omega$. Wówczas dla każdego podobiektu $C\hookrightarrow X$
\[
 C\leq\forall_K R\iff C\times K\leq R.
\]
Istotnie, faktoryzacja przez equalizator oznacza równość obu transpozycji po ograniczeniu do $C$; z bijekcji wykładniczej oznacza to, że $r$ na $C\times K$ jest stale prawdziwe, czyli że włączenie $C\times K\hookrightarrow X\times K$ faktoryzuje się przez $R$. Morfizm $m_C\times\id_K:C\times K\to X\times K$ jest monomorfizmem: równość złożeń z tym morfizmem daje równość pierwszych składowych z monomorficzności $m_C$ i równość drugich składowych wprost, a więc równość morfizmów do $C\times K$. Dla $g:Y\to X$ morfizm charakterystyczny przeciągniętego $R'=(g\times\id_K)^*R$ to $r(g\times\id_K)$. Jego transpozycją wykładniczą jest $\widetilde r g$, gdyż $\ev((\widetilde r g)\times\id_K)=r(g\times\id_K)$. Morfizm stały $X\to\Omega^K$ po złożeniu z $g$ jest odpowiednim morfizmem stałym na $Y$. Przeciągnięcie equalizatora dwóch morfizmów $u,v:X\to Z$ przez $g$ jest equalizatorem $ug,vg$: dla $h:T\to Y$ obie faktoryzacje są równoważne równości $ugh=vgh$ i jednoznaczne. Otrzymujemy $g^*(\forall_K R)=\forall_K((g\times\id_K)^*R)$.

Niech $\pi:X\times\Omega\to X$ będzie rzutem, a $Q$ podobiektem $X\times\Omega$ klasyfikowanym przez drugi rzut do $\Omega$. Kładziemy $0_X=\forall_\Omega Q$. Dla dowolnego podobiektu $B\hookrightarrow X$ o morfizmie charakterystycznym $b$ mamy $0_X\times\Omega\leq Q$. Przeciągając tę inkluzję przez $(\id_X,b):X\to X\times\Omega$, otrzymujemy $0_X\leq B$, gdyż przeciągnięcie $Q$ jest klasyfikowane przez $b$, a przeciągnięcie $0_X\times\Omega$ jest $0_X$. Zatem $0_X$ jest najmniejszy. Jest stabilny przy przeciąganiu z konstrukcji $\forall_\Omega$ i $Q$.

Dla podobiektów $A,B$ obiektu $X$ kładziemy
\[
 A\vee B=\forall_\Omega
 \bigl(((\pi^*A\Rightarrow Q)\wedge(\pi^*B\Rightarrow Q))\Rightarrow Q\bigr).
\]
Oznaczmy wyrażenie pod kwantyfikatorem przez $R$. Oznaczmy $L=(\pi^*A\Rightarrow Q)\wedge(\pi^*B\Rightarrow Q)$. Ponieważ $L\leq(\pi^*A\Rightarrow Q)$, mamy
\[
 \pi^*A\wedge L\leq \pi^*A\wedge(\pi^*A\Rightarrow Q)\leq Q.
\]
Z własności implikacji wynika $\pi^*A\leq L\Rightarrow Q=R$. Dla $B$ korzystamy z $L\leq(\pi^*B\Rightarrow Q)$ i otrzymujemy $\pi^*B\wedge L\leq Q$, a więc $\pi^*B\leq R$. Ponieważ $\pi^*A=A\times\Omega$ i $\pi^*B=B\times\Omega$ jako podobiekty $X\times\Omega$, własność $\forall_\Omega$ daje $A,B\leq A\vee B$. Jeżeli $A,B\leq D$ i $d:X\to\Omega$ klasyfikuje $D$, to przeciągnięcie $(A\vee B)\times\Omega\leq R$ przez $(\id_X,d)$ daje
\[
 A\vee B\leq ((A\Rightarrow D)\wedge(B\Rightarrow D))\Rightarrow D=D.
\]
Ostatnia równość wynika z $A,B\leq D$, co daje $A\Rightarrow D=B\Rightarrow D=1_X$, i z $1_X\Rightarrow D=D$. Zatem skonstruowany podobiekt jest najmniejszym wspólnym ograniczeniem górnym. Stabilność sum wynika ze stabilności wszystkich użytych operacji, w tym $\forall_\Omega$. Otrzymaliśmy algebrę Heytinga oraz wszystkie zadeklarowane własności przeciągania.
\end{proof}

\subsection{Wartości prawdy, negacja i logika klasyczna}

\begin{definicja}[Topos boolowski]\label{cat:def:boolean-topos}
Topos elementarny nazywamy boolowskim, jeżeli dla każdego $X$ algebra Heytinga $\operatorname{Sub}_{\mathcal E}(X)$ jest algebrą Boole'a, to znaczy każdy podobiekt $A$ spełnia $A\vee\neg A=1_X$.
\end{definicja}

W każdej algebrze Heytinga $A\wedge\neg A=0$, ponieważ $\neg A=A\Rightarrow0$. Nie wynika stąd jednak $A\vee\neg A=1$. Podwójna negacja spełnia $A\leq\neg\neg A$: pierwsza równość i własność implikacji dają tę nierówność bez prawa wyłączonego środka.

\begin{przyklad}[Logika otwartych podzbiorów]\label{cat:ex:open-heyting}
W toposie $\operatorname{Sh}(X)$ podobiekty obiektu końcowego odpowiadają otwartym podzbiorom $X$, na mocy twierdzenia~\ref{cat:thm:subobject-classifier-representation} i przykładu~\ref{cat:ex:topological-truth-values}: morfizm $1\to\Omega_J$ jest wyznaczony przez wartość na $X$, gdyż pozostałe wartości są jej ograniczeniami. Otwartemu $W$ odpowiada podsnop $1$ o wartości jednoelementowej na $U\subseteq W$ i pustej w pozostałych przypadkach. Faktoryzacja tych włączeń zachodzi dokładnie dla zawierania odpowiadających zbiorów otwartych (w szczególności sprawdzamy wartość na $W$). Stąd
\[
 W\wedge V=W\cap V,\qquad W\vee V=W\cup V,\qquad
 W\Rightarrow V=\operatorname{int}_X((X\setminus W)\cup V).
\]
Wzór na implikację wynika z faktu, że dla otwartego $T$ warunek $T\cap W\subseteq V$ jest równoważny $T\subseteq(X\setminus W)\cup V$, czyli zawieraniu w jego wnętrzu. Zatem $\neg W=\operatorname{int}_X(X\setminus W)$.

Dla $X=\mathbb R$ i $W=(0,\infty)$ mamy $\neg W=(-\infty,0)$, a ich suma nie zawiera zera. Prawo wyłączonego środka dla tego podobiektu nie zachodzi, więc $\operatorname{Sh}(\mathbb R)$ nie jest toposem boolowskim. Nie oznacza to sprzeczności z klasyczną logiką, w której prowadzimy zewnętrzne dowody: chodzi o operacje w wewnętrznej algebrze podobiektów.
\end{przyklad}

W interpretacji wewnętrznej podobiekt obiektu $X$ reprezentuje predykat na $X$, jego morfizm charakterystyczny reprezentuje wartość prawdy zależną od elementu uogólnionego, a przeciąganie reprezentuje podstawienie. Przecięcie, suma i implikacja odpowiadają spójnikom logicznym; konstrukcja $\forall_K$ interpretuje kwantyfikator ogólny dla rzutu $X\times K\to X$. Zewnętrzna równość morfizmów nie jest tym samym co prawdziwość wszystkich formuł klasycznych wewnątrz toposu.

\subsection{Morfizmy geometryczne}

\begin{definicja}[Morfizm geometryczny]\label{cat:def:geometric-morphism}
Morfizmem geometrycznym $f:\mathcal E\to\mathcal F$ nazywamy sprzężenie
\[
 f^*:\mathcal F\rightleftarrows\mathcal E:f_*,
 \qquad f^*\dashv f_*,
\]
w którym lewy funktor $f^*$ zachowuje skończone granice. Funktor $f^*$ nazywamy obrazem odwrotnym, a $f_*$ obrazem prostym. Kierunek morfizmu geometrycznego jest kierunkiem funktora $f_*$.
\end{definicja}

\begin{center}
\begin{tikzpicture}[>=stealth]
 \node (e) at (0,0) {$\mathcal E$};
 \node (f) at (5,0) {$\mathcal F$};
 \draw[->] (e) to[bend left=22] node[above] {$f_*$} (f);
 \draw[->] (f) to[bend left=22] node[below] {$f^*$} (e);
\end{tikzpicture}
\end{center}

\begin{stwierdzenie}\label{cat:prop:sheaf-geometric-inclusion}
Snopifikacja $a$ i włączenie $i:\operatorname{Sh}(\mathcal C,J)\to\widehat{\mathcal C}$ wyznaczają morfizm geometryczny $\operatorname{Sh}(\mathcal C,J)\to\widehat{\mathcal C}$ z $f^*=a$ i $f_*=i$.
\end{stwierdzenie}
\begin{proof}
Obie kategorie są toposami na mocy twierdzeń~\ref{cat:thm:grothendieck-is-elementary} i~\ref{cat:cor:presheaf-topos}. Sprzężenie $a\dashv i$ daje twierdzenie~\ref{cat:thm:sheafification}, a zachowywanie skończonych granic przez $a$ daje twierdzenie~\ref{cat:thm:sheafification-left-exact}. Spełniona jest definicja morfizmu geometrycznego. Ponieważ $i$ jest w pełni wierne, jest to przykład geometrycznego zanurzenia toposu.
\end{proof}

\begin{stwierdzenie}[Składanie morfizmów geometrycznych]\label{cat:prop:geometric-composition}
Dla morfizmów geometrycznych $f:\mathcal E\to\mathcal F$ oraz $g:\mathcal F\to\mathcal G$ złożenie $gf$ ma obraz odwrotny $f^*g^*$ i obraz prosty $g_*f_*$. Funktor identycznościowy z samym sobą daje identycznościowy morfizm geometryczny.
\end{stwierdzenie}
\begin{proof}
Dla $C\in\Ob(\mathcal G)$ i $A\in\Ob(\mathcal E)$ mamy naturalne bijekcje
\[
 \Hom_{\mathcal E}(f^*g^*C,A)
 \cong\Hom_{\mathcal F}(g^*C,f_*A)
 \cong\Hom_{\mathcal G}(C,g_*f_*A).
\]
Stąd wymagane sprzężenie. Dla skończonego diagramu $D$ w $\mathcal G$ kanoniczne izomorfizmy porównawcze dają
\[
 f^*g^*(\lim D)\cong f^*(\lim g^*D)\cong\lim f^*g^*D,
\]
więc złożenie zachowuje skończone granice. Identyczność spełnia oba warunki wprost.
\end{proof}

\begin{przyklad}[Punkt przestrzeni jako punkt toposu]\label{cat:ex:topos-point}
Niech $x\in X$. Funktor źdźbła $x^*:\operatorname{Sh}(X)\to\Set$, $F\mapsto F_x$, jest lewym sprzężeniem funktora $x_*:\Set\to\operatorname{Sh}(X)$ określonego przez
\[
 (x_*S)(U)=
 \begin{cases}S,&x\in U,\\ \{*\},&x\notin U.\end{cases}
\]
Ograniczenia między otoczeniami $x$ są identycznościami, a pozostałe są jedynymi funkcjami do zbioru jednoelementowego. Dla łańcucha otwartych zbiorów $W\subseteq V\subseteq U$, jeśli $x\in W$, wszystkie trzy wartości są $S$ i ograniczenia są identycznościami. Jeżeli $x\notin W$, oba porównywane złożenia mają jednoelementową przeciwdziedzinę, więc są równe. Ograniczenia dla identyczności są identycznościami także na zbiorze jednoelementowym. Dla funkcji $q:S\to T$ definiujemy $(x_*q)_U=q$ na otoczeniach $x$, a poza nimi jedyną funkcję między zbiorami jednoelementowymi. Naturalność wynika z identyczności ograniczeń między otoczeniami oraz jednoelementowości przeciwdziedziny w pozostałych przypadkach. Zachowanie identyczności i składania sprawdzamy na każdej składowej: albo są to odpowiednie identyczności i złożenia funkcji zbiorów, albo jedyne funkcje między singletonami. W ten sposób otrzymujemy funktor $x_*$. Dla pokrycia $U$ niezawierającego $x$ każda zgodna rodzina jest jedyna. Dla $x\in U$ przynajmniej jeden element pokrycia zawiera $x$, a każde dwa takie elementy mają przecięcie zawierające $x$. Zgodność wymusza tę samą wartość $s\in S$ na wszystkich elementach zawierających $x$, podczas gdy pozostałe wartości są jedyne. Sekcja $s$ jest więc jedynym sklejeniem (jeżeli $S=\varnothing$, takich rodzin nie ma). Stąd $x_*S$ jest snopem.

Morfizm $F\to x_*S$ jest równoważny rodzinie funkcji $F(U)\to S$ dla wszystkich otwartych otoczeń $U$ punktu $x$, zgodnej z ograniczeniami. Na pozostałych otwartych zbiorach składowe są jedyne, a naturalność dla zawierania zbioru niezawierającego $x$ w otoczeniu $x$ jest automatyczna, gdyż obie funkcje mają jednoelementową przeciwdziedzinę. Własność kogranicy źdźbła daje bijekcję $\Hom(F,x_*S)\cong\Hom(F_x,S)$: morfizm $\phi:F\to x_*S$ przechodzi na funkcję $[U,s]\mapsto\phi_U(s)$, a funkcja $q:F_x\to S$ na rodzinę $\phi_U(s)=q([U,s])$ na otoczeniach $x$ i jedyne funkcje poza nimi. Zgodność rodzin z ograniczeniami gwarantuje dobre określenie na klasach, a obie konstrukcje są odwrotne. Dla $\alpha:H\to F$ funkcja odpowiadająca $\phi\alpha$ ma wartość $\phi_U(\alpha_U(s))$ na kiełku $[U,s]$, czyli jest złożeniem funkcji odpowiadającej $\phi$ z $\alpha_x$. Dla $q:S\to T$ funkcja odpowiadająca $(x_*q)\phi$ ma wartość $q(\phi_U(s))$. To dowodzi naturalności w obu argumentach.

Źdźbło obiektu końcowego jest jednoelementowe: wszystkie otoczenia mają jedyną sekcję, a każde dwa otoczenia mają wspólne otwarte otoczenie. Dla niepustej skończonej rodziny $(F_i)_{i\in I}$ kanoniczny morfizm
\[
 (\prod_{i\in I}F_i)_x\longrightarrow\prod_{i\in I}(F_i)_x,
 \qquad [U,(s_i)_i]\longmapsto([U,s_i])_i
\]
jest dobrze określony, bo zgodność sekcji iloczynu na otoczeniu daje zgodność każdej składowej. Dla zadanej rodziny kiełków wybieramy przedstawicieli $[U_i,s_i]$ i ograniczamy je do $U=\bigcap_i U_i$, które jest otwartym otoczeniem $x$ dzięki skończoności $I$. Kiełek $[U,(s_i|_U)_i]$ ma zadany obraz. Jeśli dwa kiełki iloczynu mają ten sam obraz, to dla każdego $i$ istnieje wspólne otwarte otoczenie, na którym odpowiednie składowe są równe. Przecięcie tych skończenie wielu otoczeń jest otwartym otoczeniem $x$, na którym równe są sekcje iloczynu. Otrzymaliśmy injektywność i surjektywność kanonicznego morfizmu. Pusty iloczyn jest przypadkiem obiektu końcowego. Dla equalizatora dwóch morfizmów $F\to G$ kiełek w $F_x$ o równych obrazach ma przedstawiciela, którego obrazy są równe po ograniczeniu do pewnego otoczenia $x$. Jest to przedstawiciel w equalizatorze punktowym. Równość jego kiełków jest wykrywana tak samo jak w $F_x$, ponieważ equalizator jest punktowo podzbiorem $F$. Granice snopów są punktowe (twierdzenie~\ref{cat:thm:sheaf-limits}), a granice skończone można skonstruować z iloczynów i equalizatorów. Zatem $x^*$ zachowuje skończone granice i daje morfizm geometryczny $\Set\to\operatorname{Sh}(X)$. Ogólnie punkt toposu $\mathcal E$ to morfizm geometryczny $\Set\to\mathcal E$.
\end{przyklad}
